\chapter{Validity, Ethics, and Reproducibility}
\label{chap:limitations}

The Phase~2 results describe one configured simulation campaign. They are not
population estimates and not a validated digital twin of online climate
discourse. The harvest is therefore diagnostic: it shows which patterns
appeared under the specified graph generators, personas, articles, language
model, and scoring instruments, and where those patterns are sensitive to
design choices and API availability.

\section{Internal and statistical validity}

\paragraph{One replicate.}
The parser retained one run per configuration ($N=1$). Configurations
recorded \texttt{graphRandomSeed: 42}, but only the echo-chamber and polarised
builders consumed it; ER, small-world, scale-free, and action sampling used
unseeded \texttt{Math.random}. Article cells within a configuration share the
realised graph, persona assignment, protocol, and hosted-model context, so they
are not independent campaign replicates. Inboxes and per-article histories
were cleared between articles, and post-hoc trust updates could not alter
already completed trajectories.  The reported means and rank orders are
descriptive comparisons of this harvest; they do not provide confidence
intervals, population-level significance tests, or causal estimates of a
topology or identity-mix effect.  In particular, a small \(p\)-value computed
over pooled cells would not repair the absence of independent simulation
replicates.

\emph{Mitigation.}  Results are separated by scoring mode and accompanied by
live-cell counts and dead-cell rates.  Claims are phrased as observations
\enquote{in this harvest}, and the analysis avoids interpreting event counts
as sample size.  Future work should run independently seeded graph and model
replicates, predefine topology-level contrasts, and estimate uncertainty with
a multilevel or cluster-aware analysis whose experimental unit is the
independent run.

\paragraph{Eight-hop horizon.}
The campaign stops after eight ticks/hops, a cost compression relative to the
prior linear-chain length of 30.  This horizon is neither an empirical firestorm
duration nor Debnath et al.'s protocol: their \enquote{eight} is a skip-gram
context window of eight words.  The short horizon right-censors propagation
and weakens the interpretation of \(k^{\ast}\), which is defined as the first
observed tick above the threshold with no later observed recovery.  A cascade
that appears sustained within the window at tick eight may recover at tick nine.

\emph{Mitigation.}  Tick counts are not translated into clock time, and
\(k^{\ast}\) is reported as an eight-tick simulation outcome.  Future work
should conduct a horizon sensitivity analysis (including the 30-step
reference), use survival or censoring-aware summaries, and calibrate a time
scale only if timestamped empirical cascades become available.

\paragraph{Model dependence and stochastic services.}
Both the agents and the auditor use \texttt{gpt-4o-mini}.  The findings may
therefore reflect this model's instruction following, safety tuning,
linguistic priors, and version at the time of execution.  A hosted model can
change without a corresponding repository change, and API sampling may not be
perfectly reproducible even when the graph seed is fixed.  High misinformation
scores for conspiracy personas are compatible with prompt obedience; they are
not evidence that a field community would behave in the same way.

\emph{Mitigation.}  The model name, configuration, graph seed, raw run
directories, parser rules, and harvest timestamp are retained.  Future work
should pin provider model snapshots where possible, record sampling and API
metadata, repeat runs across several model families and temperatures, and
report model-by-condition interactions rather than presenting one model as a
generic human proxy.

\paragraph{API failure and fallback process.}
A forensic audit of the 288 selected runs found 203,928 successful responses
and 64,270 HTTP-error responses in the client counters; 121 runs recorded at
least one error.  The client made no automatic retries.  Errors were strongly
time-clustered: 57,417 occurred among runs started in the 04:00 UTC hour.
Failed rewrites propagated the unchanged input while retaining the
\texttt{reinterpret} action; failed audits left 36,194 eligible events with a
null score.  Null scores were excluded from \MI{}, \MPR{}, and \(k^{\ast}\),
but reduced \texttt{nScored} and could cause hatching.  The malformed-response
path would instead fail open to all-correct scores (\(\MI=0\)); no such parse
warning was found in the selected archived log sections.  Raw responses and
request identifiers were not retained, so provider-side uniqueness and exact
event linkage cannot be reconstructed.

The same-mode He--H sensitivity is consequential.  The canonical continuous
delta is \(+0.1549\), but becomes \(-0.0295\) when restricted to runs with
zero recorded HTTP errors.  The dual-discrete delta remains positive
(\(+0.3986\) canonical; \(+0.4695\) in zero-error runs).  This is a
conditioned sensitivity analysis, not a corrected causal estimate, because
API availability was associated with time and condition.  The continuous
collapsed H--He direction is therefore not robust to this restriction, while
the dual-discrete direction is stable across the audited checks.  Complete
forensic tables are retained separately; raw run results were not rewritten.

\emph{Mitigation.}  These failures are reported as an outcome-relevant
execution process rather than folded into cascade death.  Future runs should
use bounded backoff, fail-closed schema validation, per-request identifiers,
stage and article labels, immutable raw-response hashes, and explicit
fallback flags on events.  Propagation sparsity, rewrite failure, audit
failure, and parse failure should be separate denominators.

\paragraph{Auditor circularity.}
The generator and judge use the same \texttt{gpt-4o-mini} endpoint and client
defaults, so dual scoring tests two prompts of one model rather than two
independent instruments. High conspiracy-persona scores may primarily measure
role-prompt compliance. The campaign has no human or cross-family criterion
for accuracy, realism, or calibration. The same model family generates the
messages and judges them against five author-curated questions.  Shared model priors can create correlated generator
and judge errors, while the selected questions determine what counts as
missing or incorrect.  Discrete and continuous scores are also different
instruments: the dual-run discrete headline, its continuous sidecar, and the
continuous-only campaign are not interchangeable.  Their disagreement shows
that misinformation prevalence is partly measurement-dependent.

\emph{Mitigation.}  The two headlines are never averaged, dual gap is treated
as an instrument diagnostic rather than a third \MPR{}, and the question key
is retained for inspection.  Future validation should blind independent human
coders to condition, measure inter-rater reliability, adjudicate disagreements,
test alternative question sets, and use judge models from model families not
used to generate the content.

\paragraph{Hatched-cell policy.}
Of 1,728 thesis cells, 292 are hatched because
\texttt{nScored}$\leq 1$ (232 with zero scored events and 60 with one).
Positive configuration-level LLM usage does not prove that the article cell
itself received a successful auditor call.  They are kept in the ledger
but excluded from live-only means; they are neither zeros nor evidence that a
mix prevented misinformation.  If death depends on topology, persona, or
scoring mode, complete-case means describe surviving cascades and may be
selection-biased.  This concern is material for slices such as hierarchical
and polarised heterogeneous runs, which have comparatively high dead rates.

\emph{Mitigation.}  Every live mean is paired with its denominator and hatch
rate, and no dead cell is imputed as zero.  Future work should diagnose the
mechanism producing unscored cells, predefine a minimum-information rule,
rerun failed cells under the same protocol, and report sensitivity bounds or
selection models rather than relying only on complete cases.

\section{Construct and external validity}

\paragraph{Descriptive comparisons, not interventions.}
The grid changes several consequential features at once.  Homogeneous and
heterogeneous labels encode different persona compositions, while generators
also differ in degree sequence, reachability, and cluster assignment.
Collapsed H--He differences therefore do not isolate \enquote{diversity}, and
topology rank differences do not identify a causal network effect.  The
pooled tables concatenate cells; they do not form a physical super-graph.

\emph{Mitigation.}  Same-mode contrasts are kept separate, composition is
shown by persona family and mix, and D-net is reported outside the eight-node
pool.  Future experiments should use matched compositions, graph
randomisations that preserve relevant degree sequences, factorial ablations,
and preregistered estimands for identity, topology, and their interaction.

\paragraph{Persona reduction.}
The personas are theory-faithful prompt reductions of discourse types, not
individuals, survey respondents, or HDBSCAN centroids reconstructed from the
full Debnath corpus.  Rich and changing identities are compressed into short,
stable instructions.  The simulation therefore risks stereotyping
conspiracy-oriented, activist, scientific, or environmental participants and
may overstate within-group consistency.  A model's performance of a persona
must not be interpreted as a psychological diagnosis or a moral ordering of
real people.

\emph{Mitigation.}  Results are attributed to prompt configurations rather
than demographic groups, and the auditor's scientific key evaluates claims
rather than the worth of a persona.  Future work should co-design personas
with domain experts and affected communities, represent uncertainty and
within-group variation, test multiple prompt formulations, and audit outputs
for caricature and disparate treatment.

\paragraph{Hashtag graph, hydration failure, and absent empirical MPR.}
The 63-node D-net object is a hashtag co-occurrence reconstruction with 228
directed importer edges; it is explicitly marked
\texttt{notARetweetCascade=true}.  The edge container must not be read as
observed retweets, replies, or follower ties.  The study did not hydrate the
reported 814,924 tweet identifiers. The accessible OSF identifiers were stored
in lossy scientific notation, the Mendeley identifier dump was unavailable to
the study, and the researchers had no Twitter/X bearer token.  The fallback therefore
cannot reproduce user-level diffusion or temporal ordering.  Debnath et al.
did not report the present 0--5 misinformation index or \MPR{}, so there is no
empirical Twitter \MPR{} against which simulated \MPR{} can be matched.

\emph{Mitigation.}  The graph is labelled as a structural hashtag fallback,
tweet IDs are not invented, and simulated auditor scores are never described
as observed Twitter misinformation.  Future work requires a lawful,
lossless identifier source and approved API access, followed by an explicit
separation of hashtag co-occurrence, retweet, reply, mention, and follower
networks.  Any empirical content outcome would also require a separately
validated annotation protocol rather than retrofitting simulated \MPR{}.

\paragraph{Distributional validation with \(n_{\mathrm{real}}=1\).}
The KS and Jensen--Shannon comparisons use one empirical graph against
simulated cascades.  With \(n_{\mathrm{real}}=1\), a two-sample KS result has
little power and does not estimate natural variation among empirical
firestorms.  The DTFS of \(0.2754\) falls below the configured validation threshold of
\(0.70\). Individual metric matches do not alter this failed implementation
diagnostic.  Breadth similarity alone does not establish a
digital twin.

\emph{Mitigation.}  These quantities are presented as diagnostics and the
failed threshold is retained.  Future work should collect multiple comparable
empirical cascades, define metrics and tolerances before looking at simulation
outputs, use uncertainty-aware posterior predictive checks, and validate
structure, timing, content, and intervention response separately.

\section{Reproducibility and resource constraints}

The repository retains configs, model names, seeds, raw outputs, parser code,
derived tables, hatch rules, and the official harvest timestamp.  These
artefacts support computational traceability: a reader can follow a reported
number to a cell and inspect the transformation.  They do not guarantee exact
reproduction because the external model service, model weights, API defaults,
and provider-side moderation can change.  Re-running the full campaign also
incurs substantial monetary, energy, and time costs, and indiscriminate reruns
could overwrite or silently diverge from the recorded harvest.

\emph{Mitigation.}  Phase~2 outputs are isolated from Phase~1, completed runs
are preserved, secrets are excluded, and analysis can be regenerated without
new LLM calls.  A reproducibility release should add immutable checksums,
software and package lockfiles, model snapshot identifiers where available,
machine-readable provenance, exact cost/token totals, and a small deterministic
fixture for testing parsers without paid calls.  Replication plans should
budget API calls in advance and prioritise statistically informative
replicates over a larger but unreplicated grid.

\section{Privacy, research ethics, and dual use}

\paragraph{API privacy and cost.}
Sending content to hosted APIs transfers data to a third party and may create
retention, jurisdiction, and terms-of-service obligations.  Although the present campaign uses curated articles and synthetic personas, a
future empirical collection could process usernames, tweet text, profile
fields, or subsequently deleted content.  Public
availability does not remove contextual privacy or justify redistributing
personal data.  API access is also financially unequal and can make nominally
reproducible research inaccessible to groups without equivalent budgets.

\emph{Mitigation.}  The failed hydration workflow discarded sampled tweet text
and user fields after aggregate hashtag extraction, did not fabricate IDs,
and did not expose credentials.  Any future empirical collection should
undergo institutional ethics review, follow platform terms and data-minimisation
principles, store identifiers separately with access controls, honour deletion
where feasible, publish aggregates instead of raw personal data, document
retention, and disclose expected and realised API costs.

\paragraph{Dual-use risk.}
A simulator that identifies combinations of topology, identity prompts, and
message framing associated with high propagation or distortion could support
defensive stress testing, but it could also be used to optimise propaganda,
harassment, astroturfing, or evasion of moderation.  Detailed persona prompts
and high-performing cascade recipes can lower the cost of such misuse.

\emph{Mitigation.}  The thesis should report scientific findings without
packaging operational targeting instructions, avoid claims about vulnerable
real-world groups, and release sensitive artefacts proportionately.  Future
work should include a formal dual-use review, threat modelling, access tiers
for potentially actionable configurations, abuse monitoring, and evaluation
of defensive applications such as fact-check timing without deploying
content to live platforms.

\paragraph{No human evaluation.}
Human-evaluation templates exist in run directories, but no completed,
ethics-approved human study was conducted. There is therefore no evidence
that people find the rewrites realistic, that experts agree with the auditor,
or that the personas faithfully represent lived discourse. The absence of
human evaluation is a criterion-validity limitation as well as an ethical
safeguard: incomplete templates must not be presented as participant data.
The computational campaign used curated public texts and synthetic personas.
No formal ethics-determination, review, exemption, or approval artefact for
this master's thesis is present in the repository.
This sentence records that absence.
It does not claim that review was not required, that an exemption was
granted, or that the work has been institutionally cleared.
\UserInput{confirm with the responsible TUM school or ethics body whether
review, exemption, or a not-required determination applies, and deposit the
written artefact}
Licensing of OSF/Mendeley identifier access remains a submission-time
documentation requirement.

\emph{Mitigation and future work.}  A follow-up study should obtain ethics
approval before recruitment, use informed consent and fair compensation,
minimise exposure to harmful misinformation, provide withdrawal and debriefing
procedures, and recruit both subject-matter experts and lay raters for distinct
questions.  The protocol should preregister realism and factuality rubrics,
blind raters to condition, measure inter-rater agreement, and compare human,
independent-model, and current-auditor scores.  Human review should complement,
not merely rubber-stamp, automated scoring.

\section{Priority future work}

The next study should prioritise a smaller, preregistered validation design
over a larger unreplicated grid. Such a study should use independent seeds, matched
compositions, longer horizons, cross-model generation and judging, and blinded
human annotation.  In parallel, empirical comparison should remain structural
until multiple lawfully obtained, timestamped cascades and a validated content
measure are available.  Only after those conditions are met would it be
reasonable to test claims about generalisable topology effects, real-world
firestorm timing, or predictive correspondence with social-media
misinformation.
